```latex
\section{Wykonane zadania}

W ramach prac nad projektem listy dwukierunkowej uporządkowano
strukturę kodu źródłowego C++. Kod został rozdzielony na pliki
nagłówkowe z rozszerzeniem \texttt{.h} oraz pliki
implementacyjne z rozszerzeniem \texttt{.cpp}.

Podział ten ma na celu oddzielenie deklaracji klas, struktur
i metod od ich implementacji. Dzięki temu kod jest bardziej
czytelny, łatwiejszy w utrzymaniu i dalszym rozwijaniu.
Ułatwia również odnajdywanie poszczególnych elementów programu
oraz wprowadzanie kolejnych zmian.

Zaktualizowano pliki projektu, uwzględniając wprowadzony podział
i organizację kodu. Przygotowane zmiany przesłano do repozytorium
GitHub, dzięki czemu aktualna wersja plików jest dostępna
w repozytorium i może być dalej rozwijana.

Prace objęły również przygotowanie dokumentacji projektu
i konfigurację narzędzia Doxygen, które służy do generowania
dokumentacji na podstawie komentarzy umieszczonych bezpośrednio
w kodzie źródłowym.

\section{Niewykonane zadania}

Do dalszej realizacji pozostaje sprawdzenie kompletności
komentarzy dokumentacyjnych Doxygen w plikach źródłowych,
weryfikacja konfiguracji pliku \texttt{Doxyfile} oraz
sprawdzenie wygenerowanej dokumentacji HTML.

\section{Napotkane problemy}

Podczas prac nad projektem konieczne było uporządkowanie
konfiguracji dokumentacji Doxygen, w szczególności wskazanie
plików źródłowych C++ przeznaczonych do analizy.
Dodatkowo pojawiły się problemy z kompilacją dokumentu
raportu w LaTeX, które wymagają sprawdzenia konfiguracji
i poprawności poleceń.

\section{Zadania na kolejny tydzień}

\begin{itemize}
    \item Sprawdzenie poprawności podziału kodu na pliki
    nagłówkowe i implementacyjne.
    \item Weryfikacja aktualności plików w repozytorium GitHub.
    \item Uzupełnienie komentarzy Doxygen w kodzie źródłowym.
    \item Poprawienie konfiguracji pliku \texttt{Doxyfile}.
    \item Wygenerowanie i sprawdzenie dokumentacji HTML.
    \item Dalsze sprawdzanie działania listy dwukierunkowej
    i jej dostępnych metod.
\end{itemize}

\section{Stan dokumentacji projektowej}

Dokumentacja projektowa jest opracowywana równolegle
z porządkowaniem kodu źródłowego. Przygotowano strukturę
raportu oraz opisano organizację projektu. Rozdzielenie
plików nagłówkowych i implementacyjnych ułatwia przedstawienie
budowy programu i jego poszczególnych elementów.

Rozpoczęto również prace związane z dokumentacją Doxygen,
która ma umożliwić wygenerowanie dokumentacji technicznej
bezpośrednio na podstawie kodu źródłowego i komentarzy
dokumentacyjnych.
```